\begin{lemma}
Let \(i<k\).  If the rank of the second coordinate of \(v_i\) over the prefix
\((v_0,\ldots,v_{i-1})\) is at most \(t\), then
\[
  \psi_v(i+1)=\ell_{r}^{(v_0,\ldots,v_{i-1})},
\]
where \(r\) is that rank.
\end{lemma}
\begin{proof}
This is the defining in-range case of
\(\noderef{StarProductConcreteUnmarkedPsi}\).  The hypotheses \(i<k\) give
\(1\le i+1\le k\), so the definition uses the prefix of length \(i\) and the
child \(v_i\).  The assumed rank bound \(r\le t\) rules out the fallback value
\(0\), and the definition therefore returns the layer choice
\(\ell_r^{(v_0,\ldots,v_{i-1})}\).
\end{proof}
